% !TeX program = lualatex
% !TeX encoding = utf8
% !TeX spellcheck = uk_UA
% !BIB program = biber
% !TeX root = ../test.tex


%% ==========================================================
%% Вставити після рядка 541 (після Attention "...дротах, так само, як вода тече по трубі.")
%% і перед \section{Скін-ефект} (рядок 544) у файлі MaxwellEqns.tex
%%
%% ПЛАН МАЛЮНКІВ (усі --- заглушки, tikz будуть додані пізніше):
%%   1. pic:FieldLinesTension  --- пружні нитки: натяг уздовж + бічне відштовхування
%%   2. pic:StressElement      --- довільний майданчик dS, нормаль n, сила dF = T n dS
%%   3. pic:StressDecoding     --- однорідне поле вздовж x: майданчик n=e_x (натяг) і n=e_y (тиск)
%%   4. pic:WireInField        --- провідник у зовнішньому однорідному полі (згущення/розрідження ліній)
%%   5. pic:PinchEffect        --- плазмовий шнур, колові лінії, тиск до осі
%%   6. pic:ChargedDroplet     --- заряджена крапля, радіальні лінії, тиск назовні
%% ==========================================================

%% --------------------------------------------------------
\section{Натяги Максвелла: тиск і натяг силових ліній}\label{sec:MaxwellStresses}
%% --------------------------------------------------------

\subsection*{Від потоку енергії до механічної сили}

Теорема Пойнтінга дала нам кількісний опис того, \emphx{як} енергія електромагнітного поля рухається в просторі: вектор $\vect\Pi$ вказує
напрям і величину потоку енергії крізь одиничну площадку за одиницю часу. Але енергія не може передаватися без сили --- це фундаментальний
факт механіки. Якщо поле \emphx{переносить} енергію крізь поверхню, воно неминуче \emphx{чинить механічну дію} на цю поверхню. І справді,
розв'язуючи задачі про конденсатор і соленоїд енергетичним методом (гл.~\ref{ElEnergy}, гл.~\ref{ElMagInduction}), ми вже переконалися: поле
діє на провідники й діелектрики з поверхневою силою, яка за модулем дорівнює густині енергії. Тиск, що розпирає обмотку соленоїда, і натяг,
що стягує обкладки конденсатора, --- обидва чисельно рівні $w$.

Чи це випадковий збіг для двох окремо взятих геометрій? Виявляється, ні --- це \emphx{загальна властивість} електромагнітного поля.
Механічні натяги, що виникають у будь-якій задачі про пондеромоторні сили, завжди мають таку величину, ніби силові лінії поводяться як
натягнуті пружні нитки, які прагнуть скоротитися (натяг уздовж лінії) і взаємно відштовхуються (бічний тиск упоперек ліній). Причому і натяг,
і тиск за модулем дорівнюють об'ємній густині енергії поля. Це уявлення, запроваджене Фарадеєм і Максвеллом, дає змогу \emphx{визначати
напрям і характер механічних сил у полі, не обчислюючи їх явно} --- досить подивитися на орієнтацію силових ліній відносно поверхні.

\begin{figure}[h!]\centering
	\localinput{field-lines-tension.tikz}
	\caption{Силові лінії поля як пружні нитки: уздовж лінії --- натяг (нитка прагне скоротитися), між сусідніми паралельними
		лініями --- бічний тиск (нитки взаємно відштовхуються)\label{pic:FieldLinesTension}}
\end{figure}

\begin{Attention}
	Варто одразу застерегти від буквального розуміння цієї картини: \emphx{жодних фізичних \enquote{ниток} у просторі не існує} --- силові лінії
	є суто геометричним образом, який ми вводимо для наочного зображення поля. Проте уявлення Фарадея й Максвелла залишається надзвичайно
	корисним інструментом: воно дає простий спосіб \emphx{визначити напрям і характер} механічних сил у електромагнітному полі, навіть коли пряме
	обчислення через силу Кулона чи силу Ампера було б громіздким.
\end{Attention}

\subsection*{Тензор натягів Максвелла}

Щоб описати силу на площадку довільної орієнтації, натяг і тиск зручно об'єднати в один об'єкт --- \emphx{тензор натягів Максвелла}. Для
площадки з одиничною нормаллю $\vect n$ сила, що діє на неї з боку поля (рис.~\ref{pic:StressElement}),
\begin{equation}\label{eq:StressDefinition}
	d\vect F = \hat{T}\cdot d\vect S, \qquad d\vect S = \vect n\,dS,
\end{equation}
де $\hat{T}$ --- тензор другого рангу, а крапка позначає згортку тензора з вектором за одним індексом:
\begin{equation}\label{eq:StressComponent}
	dF_i = \sum_k T_{ik}\,dS_k.
\end{equation}

Для електричного поля компоненти тензора натягів дорівнюють
\begin{equation}\label{eq:MaxwellStressTensor}
	\tcbhighmath{
		T_{ik} = \frac{\epsilon}{4\pi}\left(E_iE_k - \frac12\delta_{ik}E^2\right),
	}
\end{equation}
а для магнітного (заміною $\Efield\to\Hfield$, $\epsilon\to1/\mu$) ---
\begin{equation}\label{eq:MaxwellStressTensorB}
	\tcbhighmath{
		T_{ik} = \frac{1}{4\pi\mu}\left(H_iH_k - \frac12\delta_{ik}H^2\right).
	}
\end{equation}

У матричному вигляді для електричного поля:
\begin{equation}\label{eq:MaxwellStressTensorMatrix}
	\hat{T}_E =
	\frac{\epsilon}{4\pi}
	\begin{pmatrix}
		\dfrac{E_x^2 - E_y^2 - E_z^2}{2} & E_xE_y & E_xE_z \\[2mm]
		E_yE_x & \dfrac{E_y^2 - E_x^2 - E_z^2}{2} & E_yE_z \\[2mm]
		E_zE_x & E_zE_y & \dfrac{E_z^2 - E_x^2 - E_y^2}{2}
	\end{pmatrix},
\end{equation}
а для магнітного:
\begin{equation}\label{eq:MaxwellStressTensorBMatrix}
	\hat{T}_H =
	\frac{1}{4\pi\mu}
	\begin{pmatrix}
		\dfrac{H_x^2 - H_y^2 - H_z^2}{2} & H_xH_y & H_xH_z \\[2mm]
		H_yH_x & \dfrac{H_y^2 - H_x^2 - H_z^2}{2} & H_yH_z \\[2mm]
		H_zH_x & H_zH_y & \dfrac{H_z^2 - H_x^2 - H_y^2}{2}
	\end{pmatrix}.
\end{equation}

Слід тензора дорівнює взятій зі знаком мінус густині енергії відповідного поля:
\begin{equation}\label{eq:StressTrace}
	\operatorname{Tr}\hat{T}_E = T_{xx}+T_{yy}+T_{zz} = -\frac{\epsilon E^2}{8\pi} = -w_E,
	\qquad
	\operatorname{Tr}\hat{T}_H = -\frac{H^2}{8\pi\mu} = -w_H.
\end{equation}

\begin{figure}[h!]\centering
	\localinput{stress-element.tikz}
	\caption{Довільно орієнтований майданчик $dS$ з нормаллю $\vect n$ у полі: результуюча сила $d\vect F=\hat T\cdot d\vect S$ в
		загальному випадку не колінеарна нормалі\label{pic:StressElement}}
\end{figure}

\subsection*{Окремий випадок: поле вздовж однієї осі}

Нехай $\Efield=E\vect e_x$ (поле напрямлене вздовж осі $x$). Тоді $E_y=E_z=0$, $E^2=E^2$, і матриця~\eqref{eq:MaxwellStressTensorMatrix}
набуває діагонального вигляду:
\begin{equation}\label{eq:StressDiagonal}
	\hat{T}_E =
	\begin{pmatrix}
		+w_E & 0 & 0 \\
		0 & -w_E & 0 \\
		0 & 0 & -w_E
	\end{pmatrix},
	\qquad w_E = \frac{\epsilon E^2}{8\pi}.
\end{equation}
Це наочно демонструє обидва ефекти: компонента $T_{xx}=+w_E$ описує \emphx{натяг уздовж поля}, а компоненти $T_{yy}=T_{zz}=-w_E$ ---
\emphx{тиск упоперек поля}.

Переконаймося, що це узгоджується з уже відомим результатом (рис.~\ref{pic:StressDecoding}). Для майданчика $\vect n=\vect e_x$
(перпендикулярного полю) $d\vect S=\vect e_x\,dS$, і
\begin{equation*}
	d\vect F = \hat{T}_E\cdot d\vect S =
	\begin{pmatrix}
		+w_E & 0 & 0 \\
		0 & -w_E & 0 \\
		0 & 0 & -w_E
	\end{pmatrix}
	\begin{pmatrix}
		dS \\ 0 \\ 0
	\end{pmatrix}
	= +w_E\,dS\,\vect e_x
\end{equation*}
--- сила напрямлена вздовж поля, це натяг, який тягне обкладки конденсатора (гл.~\ref{ElEnergy}). Для майданчика $\vect n=\vect e_y$
(паралельного полю) $d\vect S=\vect e_y\,dS$, і
\begin{equation*}
	d\vect F = \hat{T}_E\cdot d\vect S =
	\begin{pmatrix}
		+w_E & 0 & 0 \\
		0 & -w_E & 0 \\
		0 & 0 & -w_E
	\end{pmatrix}
	\begin{pmatrix}
		0 \\ dS \\ 0
	\end{pmatrix}
	= -w_E\,dS\,\vect e_y
\end{equation*}
--- сила напрямлена проти нормалі, це тиск, який розпирає обмотку соленоїда (гл.~\ref{ElMagInduction}). Обидва вже знайомі результати ---
лише два окремі майданчики одного й того самого тензора, тепер~\eqref{eq:MaxwellStressTensor} одразу дає силу для будь-якої орієнтації
площадки.

\begin{figure}[h!]\centering
	\localinput{stress-decoding.tikz}
	\caption{Однорідне поле вздовж осі $x$: майданчик $\vect n=\vect e_x$ відчуває натяг (сила вздовж поля, всередину), а
		майданчик $\vect n=\vect e_y$ --- тиск (сила проти нормалі, назовні)\label{pic:StressDecoding}}
\end{figure}

\begin{Attention}
	Якщо присутні обидва поля, повний тензор натягів --- сума $\hat{T}=\hat{T}_E+\hat{T}_H$. У матричному вигляді це означає, що
	відповідні компоненти просто додаються:
	\begin{equation*}
		T_{ik} = \frac{\epsilon}{4\pi}\left(E_iE_k - \frac12\delta_{ik}E^2\right)
		+ \frac{1}{4\pi\mu}\left(H_iH_k - \frac12\delta_{ik}H^2\right).
	\end{equation*}
	Для однорідних полів, напрямлених вздовж однієї осі, діагональні компоненти повного тензора дорівнюють $+w_E+w_H$ (уздовж спільного
	напрямку) і $-(w_E+w_H)$ (упоперек).
\end{Attention}

\subsection*{Приклад: прямий провідник в однорідному магнітному полі}

Особливо наочно картина натягу й тиску працює тоді, коли треба лише \emphx{визначити напрям} сили, не обчислюючи її величину. Нехай прямий
провідник зі струмом $I$ вміщено в однорідне магнітне поле $\Hfield_0$, силові лінії якого --- паралельні прямі (рис.~\ref{pic:WireInField}).
Власне поле струму --- це, як відомо, концентричні кола навколо провідника. Накладемо обидві картини одна на одну: з одного боку від провідника
лінії власного поля течуть \emphx{в один бік} із лініями $\Hfield_0$, і сумарне поле там \emphx{згущується}; з іншого боку вони течуть
\emphx{назустріч} одна одній, частково компенсуючись, --- і сумарне поле там \emphx{розріджується}.

\begin{figure}[h!]\centering
	\localinput{wire-in-uniform-field.tikz}
	\caption{Прямий провідник зі струмом в однорідному магнітному полі: складання власного поля струму (концентричні кола) з
		однорідним полем $\Hfield_0$ дає згущення ліній з одного боку провідника й розрідження --- з іншого\label{pic:WireInField}}
\end{figure}

Але паралельні силові лінії, як ми щойно бачили, взаємно \emphx{відштовхуються впоперек} --- отже там, де вони згущені, цей бічний тиск
найбільший. Провідник виявляється затиснутим між ділянкою підвищеного тиску (де лінії згущені) і ділянкою пониженого тиску (де вони розріджені),
а тому виштовхується в бік розрідження --- \emphx{перпендикулярно} і до провідника, і до початкового поля $\Hfield_0$. Це саме той напрям сили
Ампера $\vect F = \dfrac{I}{c}\left[\vect\ell\times\Hfield_0\right]$, який ми вже отримали раніше (гл.~\ref{MagFieldVacuum}) прямим
обчисленням, --- тепер же він випливає з самої лише картини силових ліній, без жодної формули.

\subsection*{Приклад: пінч-ефект у плазмі}

Ще один приклад, де вже недостатньо простого правила \enquote{тиск чи натяг} --- потрібна саме тензорна картина орієнтації ліній відносно
поверхні. Нехай циліндричний плазмовий шнур радіуса $a$ несе повний струм $I$ уздовж своєї осі. Магнітне поле $H(r)$, яке цей струм створює,
охоплює шнур концентричними колами --- на відміну від соленоїда, тут силові лінії \emphx{замкнені навколо} струму, а не паралельні йому
(рис.~\ref{pic:PinchEffect}). Уздовж кожної такої колової лінії поле чинить натяг (лінія \enquote{хоче} стягнутися в кільце меншого радіуса), а
це означає \emphx{тиск, спрямований до осі шнура}: поле стискає плазму, немов обруч.

\begin{figure}[h!]\centering
	\localinput{pinch-effect.tikz}
	\caption{Пінч-ефект: колові силові лінії магнітного поля навколо струму в плазмовому шнурі створюють тиск, спрямований
		до осі шнура\label{pic:PinchEffect}}
\end{figure}

Це --- \emphx{пінч-ефект} (від англ.\ \emph{pinch}, \enquote{защемлення}): струмонесучий плазмовий шнур стискається власним магнітним полем.
Стиснення триває, доки магнітний тиск не зрівноважиться тепловим (кінетичним) тиском самої плазми:
\begin{equation}\label{eq:PinchBalance}
	\frac{H^2}{8\pi} \sim n k_{\text{B}} T,
\end{equation}
де $n$ --- концентрація частинок плазми, $T$ --- її температура. Саме пінч-ефект лежить в основі найпростіших схем магнітного утримання плазми в
термоядерних установках.

\subsection*{Приклад: заряджена крапля}

Натяг і тиск силових ліній можна побачити й на викривленій поверхні. Уявіть заряджену провідну краплю (наприклад, краплю ртуті чи мильну
бульбашку з надлишковим зарядом): силові лінії виходять з її поверхні перпендикулярно, назовні, в усіх напрямках (рис.~\ref{pic:ChargedDroplet}).
Як і в соленоїді, тут поле чинить \emphx{тиск упоперек ліній}, а лінії йдуть уздовж нормалі до поверхні --- отже цей тиск $p_E=\epsilon E^2/8\pi$
штовхає поверхню назовні, \enquote{роздуваючи} краплю.

\begin{figure}[h!]\centering
	\localinput{charged-droplet.tikz}
	\caption{Заряджена крапля: радіальні силові лінії створюють тиск, спрямований назовні, якому протидіє поверхневий
		натяг рідини\label{pic:ChargedDroplet}}
\end{figure}

Цьому тиску протидіє поверхневий натяг рідини $\sigma$, що намагається стягнути краплю в кулю мінімальної площі. Поки заряд малий, поверхневий
натяг перемагає; але коли заряд перевищує критичне значення, електричний тиск перевищує поверхневий натяг, і крапля втрачає стійкість --- вона
витягується в конус і \enquote{вистрілює} тонкими зарядженими струменями. Це явище називають \emphx{нестійкістю Релея}, і воно ж лежить в основі
електроспрею й деяких методів мас-спектрометрії.

\begin{Attention}
	В усіх чотирьох прикладах --- прямому провіднику, пінч-ефекті, конденсаторі й соленоїді, зарядженій краплі --- ми жодного разу не
	підсумовували сили Кулона чи Ампера по окремих зарядах і струмах: досить було визначити, \emphx{як орієнтовані силові лінії відносно
	поверхні}. Це та сама механічна дія поля, що й потік енергії через вектор Пойнтінга, --- лише інша, \enquote{силова}, а не
	\enquote{енергетична}, проєкція одного й того самого факту: поле не просто переносить енергію, воно ще й чинить цілком реальну механічну дію
	на все, крізь що проходить. Коли пізніше ми дійдемо до електромагнітних хвиль, побачимо, що й тиск хвилі при її поглинанні речовиною --- не
	окремий ефект, а той самий натяг поля, який хвиля \enquote{несе} з собою.
\end{Attention}